% Options for packages loaded elsewhere
% Options for packages loaded elsewhere
\PassOptionsToPackage{unicode}{hyperref}
\PassOptionsToPackage{hyphens}{url}
\PassOptionsToPackage{dvipsnames,svgnames,x11names}{xcolor}
%
\documentclass[
  letterpaper,
  DIV=11,
  numbers=noendperiod]{scrreprt}
\usepackage{xcolor}
\usepackage{amsmath,amssymb}
\setcounter{secnumdepth}{5}
\usepackage{iftex}
\ifPDFTeX
  \usepackage[T1]{fontenc}
  \usepackage[utf8]{inputenc}
  \usepackage{textcomp} % provide euro and other symbols
\else % if luatex or xetex
  \usepackage{unicode-math} % this also loads fontspec
  \defaultfontfeatures{Scale=MatchLowercase}
  \defaultfontfeatures[\rmfamily]{Ligatures=TeX,Scale=1}
\fi
\usepackage{lmodern}
\ifPDFTeX\else
  % xetex/luatex font selection
\fi
% Use upquote if available, for straight quotes in verbatim environments
\IfFileExists{upquote.sty}{\usepackage{upquote}}{}
\IfFileExists{microtype.sty}{% use microtype if available
  \usepackage[]{microtype}
  \UseMicrotypeSet[protrusion]{basicmath} % disable protrusion for tt fonts
}{}
\makeatletter
\@ifundefined{KOMAClassName}{% if non-KOMA class
  \IfFileExists{parskip.sty}{%
    \usepackage{parskip}
  }{% else
    \setlength{\parindent}{0pt}
    \setlength{\parskip}{6pt plus 2pt minus 1pt}}
}{% if KOMA class
  \KOMAoptions{parskip=half}}
\makeatother
% Make \paragraph and \subparagraph free-standing
\makeatletter
\ifx\paragraph\undefined\else
  \let\oldparagraph\paragraph
  \renewcommand{\paragraph}{
    \@ifstar
      \xxxParagraphStar
      \xxxParagraphNoStar
  }
  \newcommand{\xxxParagraphStar}[1]{\oldparagraph*{#1}\mbox{}}
  \newcommand{\xxxParagraphNoStar}[1]{\oldparagraph{#1}\mbox{}}
\fi
\ifx\subparagraph\undefined\else
  \let\oldsubparagraph\subparagraph
  \renewcommand{\subparagraph}{
    \@ifstar
      \xxxSubParagraphStar
      \xxxSubParagraphNoStar
  }
  \newcommand{\xxxSubParagraphStar}[1]{\oldsubparagraph*{#1}\mbox{}}
  \newcommand{\xxxSubParagraphNoStar}[1]{\oldsubparagraph{#1}\mbox{}}
\fi
\makeatother

\usepackage{color}
\usepackage{fancyvrb}
\newcommand{\VerbBar}{|}
\newcommand{\VERB}{\Verb[commandchars=\\\{\}]}
\DefineVerbatimEnvironment{Highlighting}{Verbatim}{commandchars=\\\{\}}
% Add ',fontsize=\small' for more characters per line
\usepackage{framed}
\definecolor{shadecolor}{RGB}{241,243,245}
\newenvironment{Shaded}{\begin{snugshade}}{\end{snugshade}}
\newcommand{\AlertTok}[1]{\textcolor[rgb]{0.68,0.00,0.00}{#1}}
\newcommand{\AnnotationTok}[1]{\textcolor[rgb]{0.37,0.37,0.37}{#1}}
\newcommand{\AttributeTok}[1]{\textcolor[rgb]{0.40,0.45,0.13}{#1}}
\newcommand{\BaseNTok}[1]{\textcolor[rgb]{0.68,0.00,0.00}{#1}}
\newcommand{\BuiltInTok}[1]{\textcolor[rgb]{0.00,0.23,0.31}{#1}}
\newcommand{\CharTok}[1]{\textcolor[rgb]{0.13,0.47,0.30}{#1}}
\newcommand{\CommentTok}[1]{\textcolor[rgb]{0.37,0.37,0.37}{#1}}
\newcommand{\CommentVarTok}[1]{\textcolor[rgb]{0.37,0.37,0.37}{\textit{#1}}}
\newcommand{\ConstantTok}[1]{\textcolor[rgb]{0.56,0.35,0.01}{#1}}
\newcommand{\ControlFlowTok}[1]{\textcolor[rgb]{0.00,0.23,0.31}{\textbf{#1}}}
\newcommand{\DataTypeTok}[1]{\textcolor[rgb]{0.68,0.00,0.00}{#1}}
\newcommand{\DecValTok}[1]{\textcolor[rgb]{0.68,0.00,0.00}{#1}}
\newcommand{\DocumentationTok}[1]{\textcolor[rgb]{0.37,0.37,0.37}{\textit{#1}}}
\newcommand{\ErrorTok}[1]{\textcolor[rgb]{0.68,0.00,0.00}{#1}}
\newcommand{\ExtensionTok}[1]{\textcolor[rgb]{0.00,0.23,0.31}{#1}}
\newcommand{\FloatTok}[1]{\textcolor[rgb]{0.68,0.00,0.00}{#1}}
\newcommand{\FunctionTok}[1]{\textcolor[rgb]{0.28,0.35,0.67}{#1}}
\newcommand{\ImportTok}[1]{\textcolor[rgb]{0.00,0.46,0.62}{#1}}
\newcommand{\InformationTok}[1]{\textcolor[rgb]{0.37,0.37,0.37}{#1}}
\newcommand{\KeywordTok}[1]{\textcolor[rgb]{0.00,0.23,0.31}{\textbf{#1}}}
\newcommand{\NormalTok}[1]{\textcolor[rgb]{0.00,0.23,0.31}{#1}}
\newcommand{\OperatorTok}[1]{\textcolor[rgb]{0.37,0.37,0.37}{#1}}
\newcommand{\OtherTok}[1]{\textcolor[rgb]{0.00,0.23,0.31}{#1}}
\newcommand{\PreprocessorTok}[1]{\textcolor[rgb]{0.68,0.00,0.00}{#1}}
\newcommand{\RegionMarkerTok}[1]{\textcolor[rgb]{0.00,0.23,0.31}{#1}}
\newcommand{\SpecialCharTok}[1]{\textcolor[rgb]{0.37,0.37,0.37}{#1}}
\newcommand{\SpecialStringTok}[1]{\textcolor[rgb]{0.13,0.47,0.30}{#1}}
\newcommand{\StringTok}[1]{\textcolor[rgb]{0.13,0.47,0.30}{#1}}
\newcommand{\VariableTok}[1]{\textcolor[rgb]{0.07,0.07,0.07}{#1}}
\newcommand{\VerbatimStringTok}[1]{\textcolor[rgb]{0.13,0.47,0.30}{#1}}
\newcommand{\WarningTok}[1]{\textcolor[rgb]{0.37,0.37,0.37}{\textit{#1}}}

\usepackage{longtable,booktabs,array}
\usepackage{calc} % for calculating minipage widths
% Correct order of tables after \paragraph or \subparagraph
\usepackage{etoolbox}
\makeatletter
\patchcmd\longtable{\par}{\if@noskipsec\mbox{}\fi\par}{}{}
\makeatother
% Allow footnotes in longtable head/foot
\IfFileExists{footnotehyper.sty}{\usepackage{footnotehyper}}{\usepackage{footnote}}
\makesavenoteenv{longtable}
\usepackage{graphicx}
\makeatletter
\newsavebox\pandoc@box
\newcommand*\pandocbounded[1]{% scales image to fit in text height/width
  \sbox\pandoc@box{#1}%
  \Gscale@div\@tempa{\textheight}{\dimexpr\ht\pandoc@box+\dp\pandoc@box\relax}%
  \Gscale@div\@tempb{\linewidth}{\wd\pandoc@box}%
  \ifdim\@tempb\p@<\@tempa\p@\let\@tempa\@tempb\fi% select the smaller of both
  \ifdim\@tempa\p@<\p@\scalebox{\@tempa}{\usebox\pandoc@box}%
  \else\usebox{\pandoc@box}%
  \fi%
}
% Set default figure placement to htbp
\def\fps@figure{htbp}
\makeatother





\setlength{\emergencystretch}{3em} % prevent overfull lines

\providecommand{\tightlist}{%
  \setlength{\itemsep}{0pt}\setlength{\parskip}{0pt}}



 


\KOMAoption{captions}{tableheading}
\makeatletter
\@ifpackageloaded{tcolorbox}{}{\usepackage[skins,breakable]{tcolorbox}}
\@ifpackageloaded{fontawesome5}{}{\usepackage{fontawesome5}}
\definecolor{quarto-callout-color}{HTML}{909090}
\definecolor{quarto-callout-note-color}{HTML}{0758E5}
\definecolor{quarto-callout-important-color}{HTML}{CC1914}
\definecolor{quarto-callout-warning-color}{HTML}{EB9113}
\definecolor{quarto-callout-tip-color}{HTML}{00A047}
\definecolor{quarto-callout-caution-color}{HTML}{FC5300}
\definecolor{quarto-callout-color-frame}{HTML}{acacac}
\definecolor{quarto-callout-note-color-frame}{HTML}{4582ec}
\definecolor{quarto-callout-important-color-frame}{HTML}{d9534f}
\definecolor{quarto-callout-warning-color-frame}{HTML}{f0ad4e}
\definecolor{quarto-callout-tip-color-frame}{HTML}{02b875}
\definecolor{quarto-callout-caution-color-frame}{HTML}{fd7e14}
\makeatother
\makeatletter
\@ifpackageloaded{bookmark}{}{\usepackage{bookmark}}
\makeatother
\makeatletter
\@ifpackageloaded{caption}{}{\usepackage{caption}}
\AtBeginDocument{%
\ifdefined\contentsname
  \renewcommand*\contentsname{Table of contents}
\else
  \newcommand\contentsname{Table of contents}
\fi
\ifdefined\listfigurename
  \renewcommand*\listfigurename{List of Figures}
\else
  \newcommand\listfigurename{List of Figures}
\fi
\ifdefined\listtablename
  \renewcommand*\listtablename{List of Tables}
\else
  \newcommand\listtablename{List of Tables}
\fi
\ifdefined\figurename
  \renewcommand*\figurename{Figure}
\else
  \newcommand\figurename{Figure}
\fi
\ifdefined\tablename
  \renewcommand*\tablename{Table}
\else
  \newcommand\tablename{Table}
\fi
}
\@ifpackageloaded{float}{}{\usepackage{float}}
\floatstyle{ruled}
\@ifundefined{c@chapter}{\newfloat{codelisting}{h}{lop}}{\newfloat{codelisting}{h}{lop}[chapter]}
\floatname{codelisting}{Listing}
\newcommand*\listoflistings{\listof{codelisting}{List of Listings}}
\makeatother
\makeatletter
\makeatother
\makeatletter
\@ifpackageloaded{caption}{}{\usepackage{caption}}
\@ifpackageloaded{subcaption}{}{\usepackage{subcaption}}
\makeatother
\usepackage{bookmark}
\IfFileExists{xurl.sty}{\usepackage{xurl}}{} % add URL line breaks if available
\urlstyle{same}
\hypersetup{
  pdftitle={Dose Response graphs in R},
  pdfauthor={Lizzie Wadsworth},
  colorlinks=true,
  linkcolor={blue},
  filecolor={Maroon},
  citecolor={Blue},
  urlcolor={Blue},
  pdfcreator={LaTeX via pandoc}}


\title{Dose Response graphs in R}
\author{Lizzie Wadsworth}
\date{2026-04-06}
\begin{document}
\maketitle

\renewcommand*\contentsname{Table of contents}
{
\hypersetup{linkcolor=}
\setcounter{tocdepth}{2}
\tableofcontents
}

\bookmarksetup{startatroot}

\chapter*{Introduction}\label{introduction}
\addcontentsline{toc}{chapter}{Introduction}

\markboth{Introduction}{Introduction}

This tutorial covers how to go from raw data in a \emph{Leishmania}
PrestoBlue assay to generating a dose-response curve and calculating
IC50 values.

The R code will be very similar for each plate, so we can set it up so
there's only a few things we need to change for each one.

\bookmarksetup{startatroot}

\chapter{Put the data into a useable
format.}\label{put-the-data-into-a-useable-format.}

Before we start, we can copy our raw plate results into an excel
spreadsheet and add a column header to each one - whether the column is
a Blank, Live/DMSO control, or which drug.

For example:

\begin{longtable}[]{@{}
  >{\raggedleft\arraybackslash}p{(\linewidth - 22\tabcolsep) * \real{0.1000}}
  >{\raggedleft\arraybackslash}p{(\linewidth - 22\tabcolsep) * \real{0.1250}}
  >{\raggedleft\arraybackslash}p{(\linewidth - 22\tabcolsep) * \real{0.0750}}
  >{\raggedleft\arraybackslash}p{(\linewidth - 22\tabcolsep) * \real{0.0750}}
  >{\raggedleft\arraybackslash}p{(\linewidth - 22\tabcolsep) * \real{0.0750}}
  >{\raggedleft\arraybackslash}p{(\linewidth - 22\tabcolsep) * \real{0.0750}}
  >{\raggedleft\arraybackslash}p{(\linewidth - 22\tabcolsep) * \real{0.0750}}
  >{\raggedleft\arraybackslash}p{(\linewidth - 22\tabcolsep) * \real{0.0750}}
  >{\raggedleft\arraybackslash}p{(\linewidth - 22\tabcolsep) * \real{0.0750}}
  >{\raggedleft\arraybackslash}p{(\linewidth - 22\tabcolsep) * \real{0.0750}}
  >{\raggedleft\arraybackslash}p{(\linewidth - 22\tabcolsep) * \real{0.0750}}
  >{\raggedleft\arraybackslash}p{(\linewidth - 22\tabcolsep) * \real{0.1000}}@{}}
\toprule\noalign{}
\begin{minipage}[b]{\linewidth}\raggedleft
Blank\_L
\end{minipage} & \begin{minipage}[b]{\linewidth}\raggedleft
Live.DMSO
\end{minipage} & \begin{minipage}[b]{\linewidth}\raggedleft
Amp\_1
\end{minipage} & \begin{minipage}[b]{\linewidth}\raggedleft
Amp\_2
\end{minipage} & \begin{minipage}[b]{\linewidth}\raggedleft
Amp\_3
\end{minipage} & \begin{minipage}[b]{\linewidth}\raggedleft
CA\_1
\end{minipage} & \begin{minipage}[b]{\linewidth}\raggedleft
CA\_2
\end{minipage} & \begin{minipage}[b]{\linewidth}\raggedleft
CA\_3
\end{minipage} & \begin{minipage}[b]{\linewidth}\raggedleft
CD9\_1
\end{minipage} & \begin{minipage}[b]{\linewidth}\raggedleft
CD9\_2
\end{minipage} & \begin{minipage}[b]{\linewidth}\raggedleft
CD9\_3
\end{minipage} & \begin{minipage}[b]{\linewidth}\raggedleft
Blank\_R
\end{minipage} \\
\midrule\noalign{}
\endhead
\bottomrule\noalign{}
\endlastfoot
15023 & 37260 & 16132 & 14192 & 14305 & 24011 & 23367 & 23755 & 19062 &
22955 & 21275 & 15449 \\
14698 & 31108 & 14743 & 14619 & 14177 & 26657 & 26935 & 26680 & 23329 &
24583 & 23355 & 14383 \\
14218 & 29177 & 14215 & 13606 & 12820 & 31640 & 34248 & 33480 & 27378 &
26452 & 26957 & 14320 \\
12773 & 24540 & 13179 & 13616 & 13612 & 36486 & 40484 & 38704 & 27685 &
26143 & 26186 & 13657 \\
12755 & 23874 & 26292 & 27863 & 26425 & 41934 & 45644 & 42626 & 26234 &
24607 & 25565 & 12776 \\
12091 & 22200 & 29741 & 30546 & 29964 & 45768 & 41631 & 42039 & 28155 &
26163 & 29775 & 12502 \\
\end{longtable}

These headings are:

\begin{itemize}
\tightlist
\item
  \textbf{Blanks:} Blank\_L and Blank\_R (so we can easily pick one or
  the other)
\item
  \textbf{Live.DMSO:} The Live and DMSO controls
\item
  \textbf{Amp\_1-3:} The three Amphotericin B columns
\item
  \textbf{CA1-3:} Caffeic acid 1-3 (example Drug 1)
\item
  \textbf{CD91-3:} CD9 1-3 (example Drug 2)
\end{itemize}

Once you've got this, you want to save the table as CSV (Comma Separated
Variables) - this is easier to read in R.

\begin{figure}[H]

{\centering \pandocbounded{\includegraphics[keepaspectratio]{SaveAsCSV.gif}}

}

\caption{Save As CSV}

\end{figure}%

\begin{tcolorbox}[enhanced jigsaw, title=\textcolor{quarto-callout-caution-color}{\faFire}\hspace{0.5em}{Caution}, toprule=.15mm, left=2mm, titlerule=0mm, bottomrule=.15mm, bottomtitle=1mm, colframe=quarto-callout-caution-color-frame, coltitle=black, colbacktitle=quarto-callout-caution-color!10!white, rightrule=.15mm, leftrule=.75mm, arc=.35mm, opacitybacktitle=0.6, breakable, toptitle=1mm, colback=white, opacityback=0]

A CSV file can \textbf{only have one sheet} - if you have multiple
sheets they will be deleted!

I recommend using a separate CSV file for each plate and give each one a
unique name e.g.~Plate1\_10Dec25\_compoundX\_compoundY.csv

\end{tcolorbox}

\bookmarksetup{startatroot}

\chapter{Set up Posit Cloud workspace, upload the CSV file, and create
an R
file.}\label{set-up-posit-cloud-workspace-upload-the-csv-file-and-create-an-r-file.}

Create a \textbf{free} account with
\href{https://posit.cloud/plans/free}{Posit Cloud} and sign in.
\emph{You will not need a paid account for this work}

Then:

\begin{enumerate}
\def\labelenumi{\arabic{enumi}.}
\item
  Click \texttt{"New\ Project"} -\textgreater{}
  \texttt{"New\ RStudio\ Project"} and wait for this to load
\item
  Click the Upload button (under Files) and upload your CSV file.
\item
  Click the New File icon -\textgreater{} R Script
\item
  Click the Save icon and give the file a name
\end{enumerate}

You are now set up to start coding in R.

\begin{figure}[H]

{\centering \pandocbounded{\includegraphics[keepaspectratio]{StartanRProject.gif}}

}

\caption{Set up Posit Cloud}

\end{figure}%

\bookmarksetup{startatroot}

\chapter{Read the CSV file into R}\label{read-the-csv-file-into-r}

In R, the first line of code will be to read the CSV file.

You will want to use \texttt{headers\ =\ TRUE} so that the first row is
saved as the column headers.

You will need to put in your filename where it says
\texttt{"Test\_data.csv"}

Any text after a \texttt{\#} is a \emph{comment} - a note to ourselves
to remind us what the code does.

\begin{Shaded}
\begin{Highlighting}[]
\CommentTok{\# Read the CSV file into R, and name it raw\_data}
\NormalTok{raw\_data }\OtherTok{\textless{}{-}} \FunctionTok{read.csv}\NormalTok{(}\StringTok{"Test\_data.csv"}\NormalTok{, }\AttributeTok{header =} \ConstantTok{TRUE}\NormalTok{)}
\end{Highlighting}
\end{Shaded}

You can run the code by highlighting the line and clicking ``Run''

\texttt{raw\_data} should then appear in your Environment - you can
click this to view the table.

Check your column headings are what you expect them to be and that there
aren't any extra (empty) table cells!

\begin{figure}[H]

{\centering \pandocbounded{\includegraphics[keepaspectratio]{ReadCSV.gif}}

}

\caption{Read the CSV file in R}

\end{figure}%

\bookmarksetup{startatroot}

\chapter{Calculate the control values (Blank, Live cells,
DMSO)}\label{calculate-the-control-values-blank-live-cells-dmso}

We can use the function \texttt{mean()} to calculate the average of a
set of numbers. We can use \texttt{\$} to select the column of our
dataframe, and \texttt{{[}1:3{]}} to select rows 1-3 (or any numbers).

\section{Calculating the blank}\label{calculating-the-blank}

If both rows of the blank are media only (Blank\_L and Blank\_R) you can
calculate the blank with the line:

\begin{Shaded}
\begin{Highlighting}[]
\NormalTok{Blank\_Average }\OtherTok{\textless{}{-}} \FunctionTok{mean}\NormalTok{(}\FunctionTok{mean}\NormalTok{(raw\_data}\SpecialCharTok{$}\NormalTok{Blank\_L), }\FunctionTok{mean}\NormalTok{(raw\_data}\SpecialCharTok{$}\NormalTok{Blank\_R))}
\end{Highlighting}
\end{Shaded}

If you have some of your wells as a media and drug control, we can
calculate the blank values separately for the different controls:

\begin{Shaded}
\begin{Highlighting}[]
\CommentTok{\# Media only control in rows 1{-}3 of Blank\_L}
\NormalTok{Blank\_Average }\OtherTok{\textless{}{-}} \FunctionTok{mean}\NormalTok{(raw\_data}\SpecialCharTok{$}\NormalTok{Blank\_L[}\DecValTok{1}\SpecialCharTok{:}\DecValTok{3}\NormalTok{])}

\CommentTok{\# If you have a media + drug (no cells control), e.g. in rows 4{-}6 of Blank\_L}
\NormalTok{Media\_drug1\_control }\OtherTok{\textless{}{-}} \FunctionTok{mean}\NormalTok{(raw\_data}\SpecialCharTok{$}\NormalTok{Blank\_L[}\DecValTok{4}\SpecialCharTok{:}\DecValTok{6}\NormalTok{])}
\end{Highlighting}
\end{Shaded}

\section{Calculating the Live/ DMSO
controls}\label{calculating-the-live-dmso-controls}

You can then calculate the Live and DMSO controls using these:

\begin{Shaded}
\begin{Highlighting}[]
\CommentTok{\# Calculate Live cells average  (Live.DMSO rows 1{-}3)}
\NormalTok{Live\_cells }\OtherTok{\textless{}{-}} \FunctionTok{mean}\NormalTok{(raw\_data}\SpecialCharTok{$}\NormalTok{Live.DMSO[}\DecValTok{1}\SpecialCharTok{:}\DecValTok{3}\NormalTok{]) }\SpecialCharTok{{-}}\NormalTok{ Blank\_Average}

\CommentTok{\# Calculate DMSO control average (Live.DMSO rows 4{-}6)}
\NormalTok{DMSO\_cells }\OtherTok{\textless{}{-}} \FunctionTok{mean}\NormalTok{(raw\_data}\SpecialCharTok{$}\NormalTok{Live.DMSO[}\DecValTok{4}\SpecialCharTok{:}\DecValTok{6}\NormalTok{]) }\SpecialCharTok{{-}}\NormalTok{ Blank\_Average}
\end{Highlighting}
\end{Shaded}

\bookmarksetup{startatroot}

\chapter{Define drug information, concentrations and
locations}\label{define-drug-information-concentrations-and-locations}

For each drug, you'll want to include the name, concentrations,
locations and which blank to use. Once you have this information, you
can use the provided functions to calculate the IC50 and draw the
dose-response curve.

You can copy the below code and add in your values, concentrations etc.

Keep the variables as \texttt{AmpB}, \texttt{Drug1} and \texttt{Drug2}
as we'll use these in the next step.

\begin{Shaded}
\begin{Highlighting}[]
\CommentTok{\# Data for AmpB}
\NormalTok{AmpB }\OtherTok{\textless{}{-}} \FunctionTok{list}\NormalTok{(}\AttributeTok{name =} \StringTok{"Amphotericin B"}\NormalTok{, }
          \AttributeTok{start\_conc =} \FloatTok{1.8}\NormalTok{, }
          \AttributeTok{unit =} \StringTok{"µM"}\NormalTok{,}
          \AttributeTok{dilution\_factor =} \DecValTok{3}\NormalTok{, }
          \AttributeTok{col\_start =} \DecValTok{3}\NormalTok{, }
          \AttributeTok{col\_stop =} \DecValTok{5}\NormalTok{, }
          \AttributeTok{blank =}\NormalTok{ Blank\_Average, }
          \AttributeTok{positive\_control =}\NormalTok{ Live\_cells)}

\CommentTok{\# Data for Drug 1}
\NormalTok{Drug1 }\OtherTok{\textless{}{-}} \FunctionTok{list}\NormalTok{(}\AttributeTok{name =} \StringTok{"Caffeic Acid"}\NormalTok{, }
             \AttributeTok{start\_conc =} \FloatTok{277.5}\NormalTok{, }
             \AttributeTok{unit =} \StringTok{"µM"}\NormalTok{,}
             \AttributeTok{dilution\_factor =} \DecValTok{2}\NormalTok{, }
             \AttributeTok{col\_start =} \DecValTok{6}\NormalTok{, }
             \AttributeTok{col\_stop =} \DecValTok{8}\NormalTok{, }
             \AttributeTok{blank =}\NormalTok{ Blank\_Average, }
             \AttributeTok{positive\_control =}\NormalTok{ DMSO\_cells)}

\CommentTok{\# Data for Drug 2 }
\NormalTok{Drug2 }\OtherTok{\textless{}{-}} \FunctionTok{list}\NormalTok{(}\AttributeTok{name =} \StringTok{"CD9"}\NormalTok{, }
             \AttributeTok{start\_conc =} \FloatTok{141.32}\NormalTok{, }
             \AttributeTok{unit =} \StringTok{"µM"}\NormalTok{,}
             \AttributeTok{dilution\_factor =} \DecValTok{2}\NormalTok{, }
             \AttributeTok{col\_start =} \DecValTok{9}\NormalTok{, }
             \AttributeTok{col\_stop =} \DecValTok{11}\NormalTok{, }
             \AttributeTok{blank =}\NormalTok{ Blank\_Average, }
             \AttributeTok{positive\_control =}\NormalTok{ DMSO\_cells)}
\end{Highlighting}
\end{Shaded}

See the below table for an explanation of each value.

\begin{longtable}[]{@{}
  >{\raggedright\arraybackslash}p{(\linewidth - 4\tabcolsep) * \real{0.3333}}
  >{\raggedright\arraybackslash}p{(\linewidth - 4\tabcolsep) * \real{0.3333}}
  >{\raggedright\arraybackslash}p{(\linewidth - 4\tabcolsep) * \real{0.3333}}@{}}
\toprule\noalign{}
\begin{minipage}[b]{\linewidth}\raggedright
Parameter
\end{minipage} & \begin{minipage}[b]{\linewidth}\raggedright
Example Value
\end{minipage} & \begin{minipage}[b]{\linewidth}\raggedright
Meaning
\end{minipage} \\
\midrule\noalign{}
\endhead
\bottomrule\noalign{}
\endlastfoot
\texttt{name} & ``Amphotericin B'' & The name you want to be annotated
on the figures for each drug \\
\texttt{start\_conc} & 1.8 & The highest concentration tested in the
assay \\
\texttt{unit} & ``µM'' & The units of the concentration given \\
\texttt{dilution\_factor} & 3 & The dilution factor by which the drug
concentration drops at each step, eg. 3 for a 3x dilution \\
\texttt{col\_start} & 3 & Which column the drug starts at (e.g.~column
3) \\
\texttt{col\_stop} & 5 & Which column the drug ends at (e.g.~column 5,
as it is in columns 3-5) \\
\texttt{blank} & \texttt{Blank\_Average} & Which of the calculated blank
values you want to use (e.g.~Blank\_Average or Media\_drug1\_control) \\
\texttt{positive\_control} & \texttt{Live\_cells} & Whether you are
comparing to Live Cells or the DMSO treated cells as a positive
control \\
\end{longtable}

\bookmarksetup{startatroot}

\chapter{Make your IC50 calculations and Dose Response
graphs}\label{make-your-ic50-calculations-and-dose-response-graphs}

\section{Install packages}\label{install-packages}

First you'll need to install the packages needed:

\begin{Shaded}
\begin{Highlighting}[]
\CommentTok{\# Install packages}
\FunctionTok{install.packages}\NormalTok{(}\StringTok{"ggplot2"}\NormalTok{)}
\FunctionTok{install.packages}\NormalTok{(}\StringTok{"drc"}\NormalTok{)}
\end{Highlighting}
\end{Shaded}

\section{Copy and run the below code}\label{copy-and-run-the-below-code}

Now all of the setup for your plate is done, the rest we can use
\textbf{the same code for each one.}

See the screencast below the code for how to use:

\emph{In this example, the ggplot graph doesn't calculate correctly for
AmpB, this is to do with the variability in the data as viability
decreases for the first few dilutions}

\begin{Shaded}
\begin{Highlighting}[]
\CommentTok{\# Import libraries}
\FunctionTok{library}\NormalTok{(ggplot2)}
\FunctionTok{library}\NormalTok{(drc)}

\CommentTok{\# Define functions}
\NormalTok{calculate\_dilutions }\OtherTok{\textless{}{-}} \ControlFlowTok{function}\NormalTok{(highest\_conc, dilution\_factor, }\AttributeTok{dilutions=}\DecValTok{5}\NormalTok{)\{}
\NormalTok{  result\_vector }\OtherTok{\textless{}{-}}\NormalTok{ highest\_conc }\SpecialCharTok{/}\NormalTok{ (dilution\_factor}\SpecialCharTok{\^{}}\NormalTok{(}\DecValTok{0}\SpecialCharTok{:}\NormalTok{dilutions))  }
  \FunctionTok{return}\NormalTok{(result\_vector)}
\NormalTok{\}}

\NormalTok{plot\_publication\_curve }\OtherTok{\textless{}{-}} \ControlFlowTok{function}\NormalTok{(drug, raw\_data\_matrix) \{}
  
  \CommentTok{\# Extract data points for the chosen drug, using col\_start and col\_stop to select columns (all rows)}
\NormalTok{  drug\_raw\_data }\OtherTok{\textless{}{-}}\NormalTok{ raw\_data\_matrix[, drug}\SpecialCharTok{$}\NormalTok{col\_start}\SpecialCharTok{:}\NormalTok{drug}\SpecialCharTok{$}\NormalTok{col\_stop]}
  
  \CommentTok{\# Calculate drug concentrations for each row using start concentration and dilution factor}
\NormalTok{  concs }\OtherTok{\textless{}{-}} \FunctionTok{calculate\_dilutions}\NormalTok{(drug}\SpecialCharTok{$}\NormalTok{start\_conc, drug}\SpecialCharTok{$}\NormalTok{dilution\_factor, }\AttributeTok{dilutions =} \DecValTok{5}\NormalTok{)}
  
  \CommentTok{\# Make a big table with each concentration repeated for each raw value}
\NormalTok{  drug\_data }\OtherTok{\textless{}{-}} \FunctionTok{data.frame}\NormalTok{(}
    \AttributeTok{conc =} \FunctionTok{rep}\NormalTok{(concs, }\AttributeTok{times =} \FunctionTok{ncol}\NormalTok{(drug\_raw\_data)),}
    \AttributeTok{raw\_val =} \FunctionTok{unlist}\NormalTok{(drug\_raw\_data)}
\NormalTok{  )}
  
  \CommentTok{\# Minus blank, and divide by positive control, to get \% viability}
\NormalTok{  drug\_data}\SpecialCharTok{$}\NormalTok{viability }\OtherTok{\textless{}{-}}\NormalTok{ ((drug\_data}\SpecialCharTok{$}\NormalTok{raw\_val }\SpecialCharTok{{-}}\NormalTok{ drug}\SpecialCharTok{$}\NormalTok{blank) }\SpecialCharTok{/}\NormalTok{ drug}\SpecialCharTok{$}\NormalTok{positive\_control) }\SpecialCharTok{*} \DecValTok{100}
  
  \CommentTok{\# Use drm from drc package to fit a dose{-}response curve to the data}
\NormalTok{  max\_viability }\OtherTok{\textless{}{-}} \FunctionTok{max}\NormalTok{(drug\_data}\SpecialCharTok{$}\NormalTok{viability)}
\NormalTok{  drug\_model }\OtherTok{\textless{}{-}} \FunctionTok{drm}\NormalTok{(viability }\SpecialCharTok{\textasciitilde{}}\NormalTok{ conc, }\AttributeTok{data =}\NormalTok{ drug\_data, }\AttributeTok{fct =} \FunctionTok{LL.4}\NormalTok{(}\AttributeTok{fixed =} \FunctionTok{c}\NormalTok{(}\ConstantTok{NA}\NormalTok{, }\DecValTok{0}\NormalTok{, max\_viability, }\ConstantTok{NA}\NormalTok{)))  }
  
  \CommentTok{\# print the results }
  \FunctionTok{print}\NormalTok{(}\FunctionTok{summary}\NormalTok{(drug\_model))}
  
  \CommentTok{\# calculate and print the IC50 value}
\NormalTok{  ic50\_info }\OtherTok{\textless{}{-}} \FunctionTok{ED}\NormalTok{(drug\_model, }\DecValTok{50}\NormalTok{, }\AttributeTok{type =} \StringTok{"absolute"}\NormalTok{, }\AttributeTok{display =} \ConstantTok{FALSE}\NormalTok{)}
  \FunctionTok{cat}\NormalTok{(}\StringTok{"IC50 value: "}\NormalTok{, ic50\_info[}\DecValTok{1}\NormalTok{], }\StringTok{" with standard error "}\NormalTok{, ic50\_info[}\DecValTok{2}\NormalTok{], }\StringTok{"}\SpecialCharTok{\textbackslash{}n}\StringTok{"}\NormalTok{)}
  
  \CommentTok{\# plot this dose{-}response model using Base R }
  \FunctionTok{plot}\NormalTok{(drug\_model, }\AttributeTok{broken =} \ConstantTok{TRUE}\NormalTok{, }\AttributeTok{type =} \StringTok{"all"}\NormalTok{, }\AttributeTok{xlab =} \FunctionTok{paste0}\NormalTok{(drug}\SpecialCharTok{$}\NormalTok{name, }\StringTok{" ("}\NormalTok{,drug}\SpecialCharTok{$}\NormalTok{unit,}\StringTok{")"}\NormalTok{), }\AttributeTok{ylab =} \StringTok{"Percent viability"}\NormalTok{)}
\NormalTok{  dose\_curve\_baseR }\OtherTok{\textless{}{-}} \FunctionTok{recordPlot}\NormalTok{()}
  
  \CommentTok{\# Use the model to create a set of predicted values, to use as a curve in a ggplot }
\NormalTok{  prediction\_data }\OtherTok{\textless{}{-}} \FunctionTok{expand.grid}\NormalTok{(}\AttributeTok{conc=}\FunctionTok{seq}\NormalTok{(}\AttributeTok{from=}\FunctionTok{min}\NormalTok{(concs), }\AttributeTok{to=}\FunctionTok{max}\NormalTok{(concs), }\AttributeTok{length =} \DecValTok{100}\NormalTok{))}
\NormalTok{  pm }\OtherTok{\textless{}{-}} \FunctionTok{predict}\NormalTok{(drug\_model, }\AttributeTok{newdata=}\NormalTok{prediction\_data, }\AttributeTok{interval=}\StringTok{"confidence"}\NormalTok{)}

\NormalTok{  prediction\_data}\SpecialCharTok{$}\NormalTok{p }\OtherTok{\textless{}{-}}\NormalTok{ pm[,}\DecValTok{1}\NormalTok{]}
\NormalTok{  prediction\_data}\SpecialCharTok{$}\NormalTok{pmin }\OtherTok{\textless{}{-}}\NormalTok{ pm[,}\DecValTok{2}\NormalTok{]}
\NormalTok{  prediction\_data}\SpecialCharTok{$}\NormalTok{pmax }\OtherTok{\textless{}{-}}\NormalTok{ pm[,}\DecValTok{3}\NormalTok{]}
  
\NormalTok{  xlabel }\OtherTok{\textless{}{-}} \FunctionTok{paste0}\NormalTok{(drug}\SpecialCharTok{$}\NormalTok{name, }\StringTok{" ("}\NormalTok{,drug}\SpecialCharTok{$}\NormalTok{unit,}\StringTok{")"}\NormalTok{)}
  
  \CommentTok{\# plot this with ggplot}
\NormalTok{  dose\_curve\_ggplot }\OtherTok{\textless{}{-}} \FunctionTok{ggplot}\NormalTok{(drug\_data, }\FunctionTok{aes}\NormalTok{(}\AttributeTok{x =}\NormalTok{ conc, }\AttributeTok{y =}\NormalTok{ viability)) }\SpecialCharTok{+} 
    \CommentTok{\# plot the raw values, add scatter points}
    \FunctionTok{geom\_point}\NormalTok{() }\SpecialCharTok{+} 
    \CommentTok{\# add a ribbon using the predicted values}
    \FunctionTok{geom\_ribbon}\NormalTok{(}\AttributeTok{data=}\NormalTok{prediction\_data, }\FunctionTok{aes}\NormalTok{(}\AttributeTok{x=}\NormalTok{conc, }\AttributeTok{y=}\NormalTok{p, }\AttributeTok{ymin=}\NormalTok{pmin, }\AttributeTok{ymax=}\NormalTok{pmax), }\AttributeTok{alpha =} \FloatTok{0.2}\NormalTok{) }\SpecialCharTok{+} 
    \FunctionTok{geom\_line}\NormalTok{(}\AttributeTok{data=}\NormalTok{prediction\_data, }\FunctionTok{aes}\NormalTok{(}\AttributeTok{x=}\NormalTok{conc, }\AttributeTok{y=}\NormalTok{p)) }\SpecialCharTok{+}
    \CommentTok{\# make the x axis a log scale}
    \FunctionTok{scale\_x\_log10}\NormalTok{() }\SpecialCharTok{+} 
    \CommentTok{\# x and y axis labels}
    \FunctionTok{xlab}\NormalTok{(xlabel) }\SpecialCharTok{+} \FunctionTok{ylab}\NormalTok{(}\StringTok{"Percent viability (cm)"}\NormalTok{) }\SpecialCharTok{+} 
    \FunctionTok{theme\_minimal}\NormalTok{()}
    
    
  \CommentTok{\# return the values }
  \FunctionTok{return}\NormalTok{(}\FunctionTok{list}\NormalTok{(}\AttributeTok{IC50 =}\NormalTok{ ic50\_info, }\AttributeTok{baseR =}\NormalTok{ dose\_curve\_baseR, }\AttributeTok{ggplot =}\NormalTok{ dose\_curve\_ggplot)) }
\NormalTok{\}}

\CommentTok{\# Run the functions for each graph}
\NormalTok{AmpB\_graphs }\OtherTok{\textless{}{-}} \FunctionTok{plot\_publication\_curve}\NormalTok{(AmpB, raw\_data)}
\NormalTok{AmpB\_graphs}\SpecialCharTok{$}\NormalTok{IC50}
\NormalTok{AmpB\_graphs}\SpecialCharTok{$}\NormalTok{baseR}
\NormalTok{AmpB\_graphs}\SpecialCharTok{$}\NormalTok{ggplot}

\NormalTok{Drug1\_graphs }\OtherTok{\textless{}{-}} \FunctionTok{plot\_publication\_curve}\NormalTok{(Drug1, raw\_data)}
\NormalTok{Drug1\_graphs}\SpecialCharTok{$}\NormalTok{IC50 }
\NormalTok{Drug1\_graphs}\SpecialCharTok{$}\NormalTok{baseR}
\NormalTok{Drug1\_graphs}\SpecialCharTok{$}\NormalTok{ggplot}

\NormalTok{Drug2\_graphs }\OtherTok{\textless{}{-}} \FunctionTok{plot\_publication\_curve}\NormalTok{(Drug2, raw\_data)}
\NormalTok{Drug2\_graphs}\SpecialCharTok{$}\NormalTok{IC50}
\NormalTok{Drug2\_graphs}\SpecialCharTok{$}\NormalTok{baseR}
\NormalTok{Drug2\_graphs}\SpecialCharTok{$}\NormalTok{ggplot}
\end{Highlighting}
\end{Shaded}

\section{Get the result}\label{get-the-result}

Each of the lines at the bottom will give you a different result:

\begin{itemize}
\tightlist
\item
  AmpB\_graphs\$IC50: The IC50 value and standard error for Amphotericin
  B
\item
  AmpB\_graphs\$baseR: A simple dose-response curve
\item
  AmpB\_graphs\$ggplot: A fancier dose-response curve (which doesn't
  work for AmpB in this example)
\end{itemize}

Once you've run the whole code, you can select and run one of these
lines at a time.

The IC50 values will appear in the console, and the plots will appear in
the bottom right panel. You can export the plots as PNG files.

\begin{figure}[H]

{\centering \pandocbounded{\includegraphics[keepaspectratio]{GetResults.gif}}

}

\caption{Run code and extract results}

\end{figure}%




\end{document}
